A lot of consideration was placed into the selection of the image processor. Initially, I investigated ASICs suited for image processing. I found that most active and supported imaging ASICs are manufactured for larger scopes i.e., video and AI processing, which exceeded the current power and thermal requirements of the project.
Most ASICs that fit the current scope are more dated and less-supported. For example, the ADSP-BF70x line were prime candidates at the start of the project since their attributes fit the current scope and aren't "overkill". They processed low-resolution images at a high-energy efficiency, could utilise external RAM and interfaced with CMOS senors with PPI.
Despite this, the line (as of mid-2025) is more than a decade old and has limited support. Other processors from NXP, Microchip and Texas Instruments either showed the same limited support or lacked imaging processing capabilities (most designed for power systems).
After consulting with online forums like \href{https://www.reddit.com/r/ECE/comments/1mizdoo/best_asic_for_image_processing/}{r/ECE} and \href{https://www.reddit.com/r/embedded/comments/1mfv3zn/comment/n6kd7jb/}{r/embedded}, as well as with Dr. Robert Howie regarding his students' previous star tracker project, it was decided that a general processor like the STM32 was best suited for the current application.
\nl
STM32s have been widely used in CubeSat applications over the years and have become a widely supported platform for ameteur aerospace applications. Unlike FPGAs, which can compute parallel processing, STM32s use high-level languages like C/++, have faster development cycles and have a larger community. They are typically more power efficient and can have built-in image processing and encoders.
Amongst the STM32 line, I selected the STM32N6 as the image processor over others like the H7 for the following reasons:
\begin{itemize}
    \item Supports DCMI as well as MIPI CSI-2. H7s do not support the latter.
    \item Uses an image signal processor (ISP), digital signal processor (DSP) and neural processing unit (NPU) for processing unlike a traditional software-powered CPU.
    \item Has hardware-accelerated subsampling techniques, whereas others use software.
    \item Has a higher power efficiency compared to the H7 (and can be underclocked to save more power).
    \item More recent and will be supported for longer.
\end{itemize}
The following table explores some of the key differences between the STM32 processors:
\begin{table}[H]
\centering
\begin{tabular}{|p{2cm}|p{3.75cm}|p{3.75cm}|p{3.75cm}|}
\hline
\textbf{Feature} & \textbf{STM32N6 Series} & \textbf{STM32H7 Series} & \textbf{STM32F4 Series} \\
\hline
Core & Arm Cortex-M55 + Neural ART NPU & Arm Cortex-M7 (single or dual-core M7/M4) & Arm Cortex-M4 \\
\hline
Performance & Up to 480 MHz (core), includes vector extensions & Up to 550 MHz (M7), strong DSP/float perf & Up to 180 MHz, moderate DSP capability \\
\hline
Memory (RAM/Flash) & Up to 4 MB SRAM, 4 MB Flash & Up to 1 MB SRAM, 2 MB Flash & Up to 256 KB SRAM, 2 MB Flash \\
\hline
Imaging Interfaces & DCMI, SPI, QuadSPI, JPEG Codec, Chrom-ART, DMA2D & DCMI, JPEG Codec, Chrom-ART, DMA2D & DCMI (on select models), basic DMA \\
\hline
Subsampling / Encoding & Hardware JPEG with native Y'UV support & Hardware JPEG + DMA2D for YUV & Software-based only \\
\hline
Power Efficiency & Moderate to High (based on workload + NPU efficiency) & Moderate (high-performance, higher dynamic power) & Good (lower performance = lower power) \\
\hline
Target Applications & AI vision, edge inference, intelligent sensors, imaging pipelines & High-performance imaging, audio, motor control, GUI apps & General-purpose MCU, audio, lower-end imaging \\
\hline
Availability & Early access or newly launched (as of 2025) & Mature and widely supported & Legacy but stable \\
\hline
\end{tabular}
\caption{Comparison of STM32N6, STM32H7, and STM32F4 Microcontroller Families}
\label{tab:stm32-family-comparison}
\end{table}
I selected the \href{https://au.mouser.com/ProductDetail/STMicroelectronics/STM32N657I0H3Q?qs=a2MtRaTmNOQAexZMsq8u8A%3D%3D}{STM32N657I0H3Q} as the image processor for the breakout board. I opted for this processor amongst its other counterparts (like the \href{https://au.mouser.com/ProductDetail/STMicroelectronics/STM32N657X0H3Q?qs=%252BXxaIXUDbq0JgAtfYs%252BpjA%3D%3D}{STM32N657X0H3Q}) because of a similar feature-list at a lower price and significantly lower pin count.
Although this component is designed for more complex- and intense-tasks, I will be able to evaluate which features will be utilised and which won't be, which will aid in selecting the optimal processor for the prototype. It serves as a solid middle ground between the more extensive and cheaper counterparts
\begin{table}[H]
\centering
\begin{tabular}{|l|l|p{8cm}|}
    \hline
\textbf{Segment} & \textbf{Code} & \textbf{Meaning / Explanation} \\
\hline
Manufacturer Family & STM32 & STMicroelectronics 32-bit microcontroller family \\
\hline
Series & N6 & STM32N6 series - based on Arm Cortex-M55 core with Neural ART NPU \\
\hline
Sub-Family / Feature Set & 57 & High-end variant in the STM32N6 family (e.g., more memory, interfaces, and image processing features) \\
\hline
Memory / Package Variant & I0 & Midrange flash/RAM size and medium-pin-count package \\
\hline
Temperature Grade & H & High-temperature / industrial range (typically -40°C to +85°C or higher) \\
\hline
Variant Code & 3 & Specific feature or silicon variant (usually minor revisions or option sets) \\
\hline
Qualification Grade & Q & Automotive or industrial qualification \\
\hline
\end{tabular}
\caption{Breakdown of STM32N657I0H3Q Part Number}
\label{tab:stm32n6-partnumber-n657i0}
\end{table}

Despite the current choice, it is highly recommended that if the scope of the project were to expand in the future, further research into ASICs should be completed.
Compared to modern SoCs or FPGAs with image processing capabilities, some ASICs provide a lower-cost entry point with greater efficiency in real-time image processing.
\nl
Since the selected processor has sufficient internal memory to buffer a full-resolution image, external memory is not required to support real-time image acquisition and processing.
If future revisions of the project (i.e., the prototype) require more memory to buffer the image, the selected processor is compatible with external uSDRAM like the \href{https://au.mouser.com/ProductDetail/ISSI/IS45S16320D-7TLA2?qs=Iq7RKQRjpgqlMjstEZ4V8g%3D%3D}{IS45S16320D-7TLA2}.
